\subsection{La corrección relativa es \texorpdfstring{$O(k^{-1})$}{O(1/k)} en la clase con brecha}\label{sec:B3T06}
Conservamos la notación de la Sección~\ref{sec:B3T05} y ponemos $q=1-p$.
La fórmula de Jacobi de la Sección~\ref{sec:B3T01} se puede escribir, en esta rama pequeña,
\[
 (\nabla\ell)_e=\frac{6p_e}{q_e(1-4p_e)}
 -\frac12\left[
 \frac{q_e^2+3p_e^2}{(1-4p_e)^2}A_eGA_e^T
 -\sum_{d\ni e}G_{dd}\right].
\]
La suma final contiene las dos diagonales, no las restricciones RC.
Esta expresión se sigue también de derivar el factor local y $S_{\rm gf}$:
$\partial_{p_e}\tau_e=(q_e^2+3p_e^2)/(1-4p_e)^2$,
mientras que el bloque diagonal RC de $S$ permanece idénticamente cero.
La Sección~\ref{sec:B3T05} implica $A_eGA_e^T=4+O(k^{-1})$ y
$\sum_{d\ni e}G_{dd}=2+O(k^{-1})$.
Además $p_e=O(k^{-1})$, por lo que
\[
 \boxed{\nabla_p\ell=-\mathbf1+O_\infty(k^{-1}).}
\]
El término constante es una dirección normal RC, no una corrección
efectiva a las probabilidades.

\paragraph{Respuesta relativa.}
El complemento de Schur del bloque $N$ da, usando formas invariantes
bajo el gauge,
\[
 R=T-TAG A^TT.
\]
Como $\|P^{-1}T\|_\infty\le2$, $\|A\|_\infty=4$ y
$\|A^TT\|_\infty=\max_\ell\sum_{e\in\ell}\tau_e\le2$ para $k$ grande,
\[
 \|P^{-1}R\|_\infty
 \le2(1+8\|G\|_\infty)=O_{\delta_0}(1).
\]
El vector constante $\mathbf1$ está en el espacio normal de las restricciones
RC, por lo que $R\mathbf1=0$. La Sección~\ref{sec:B3T01} y las dos estimaciones anteriores dan
\[
 \boxed{\delta=P^{-1}R(\nabla\ell+\mathbf1)=O_\infty(k^{-1}).}
\]
En particular, $1+\delta_e>0$ para $k$ suficientemente grande en esta clase.
La prueba no necesita una cota inferior adicional para $p_e$, pero sí
$p_e>0$, las vacancias y la brecha ponderada uniformes de la Sección~\ref{sec:B3T05}.

\paragraph{Cancelación de primer orden en la parte derivativa.}
La identidad exacta $\sum_{c\in r}p_c\delta_c=0$ de la Sección~\ref{sec:B3T01} implica
\[
 \sum_{c\in r}p_c\log(1+\delta_c)
 =-\frac12\sum_{c\in r}p_c\delta_c^2+O(k^{-3})=O(k^{-2}).
\]
Así se cancela el primer orden de la parte derivativa. Para el defecto
real de inserción sigue siendo necesario añadir $r_c$:
\[
 \mathbb E_p\epsilon
 =\mathbb E_p r-\tfrac12\mathbb E_p\delta^2+O(k^{-3}).
\]
La Sección~\ref{sec:macro} controla $r_c$. El resultado vale
uniformemente en la clase declarada; no se afirma que todo residual legal
de reinas pertenezca a ella.
